% Apuntes de teoría en PDF · 1.º ESO · Unidad 8: Áreas y perímetros (versión en castellano)
% Mateix estil que els apunts de Batxillerat (franja i caixes de teoria en blau, exemples en roig), però més visuals,
% curts i amb dibuixos, per a l'alumnat de 3r ESO. Es compila amb pdflatex dues vegades (per als enllaços de l'índex),
% amb logo-fentmates.pdf (_eines/logo/pdf/horitzontal-clar.pdf) al costat.
\documentclass[11pt,a4paper]{article}
\usepackage[utf8]{inputenc}
\usepackage[T1]{fontenc}
\usepackage[spanish,es-noshorthands,es-nodecimaldot,es-nolists]{babel}
\usepackage[sfdefault,book,lining]{FiraSans}
\usepackage[italic,noplusnominus,noequal]{mathastext}
\linespread{1.18}
\usepackage[a4paper,top=2.4cm,bottom=2.2cm,left=1.6cm,right=1.6cm,headheight=1.1cm,headsep=0.5cm,footskip=1cm]{geometry}
\usepackage{amsmath,amssymb,cancel}
\usepackage{icomma}
\usepackage{graphicx}
\usepackage[table]{xcolor}
\usepackage{tikz}\usetikzlibrary{calc,positioning,arrows.meta,decorations.pathreplacing,shapes.geometric}
\usepackage{fontawesome5}
\usepackage{ccicons}
\usepackage{fancyhdr}
\usepackage{enumitem}
\usepackage{array}
\usepackage{needspace}
\usepackage[most]{tcolorbox}
\usepackage{titlesec}
\usepackage[hidelinks]{hyperref}
\hypersetup{pdftitle={Apuntes de teoría. Unidad 8: Áreas y perímetros. 1.º ESO},pdfsubject={Matemáticas 1.º ESO},pdfkeywords={perímetro, área, unidades de longitud, unidades de superficie, cuadrado, rectángulo, romboide, triángulo, rombo, trapecio, polígonos regulares, circunferencia, círculo, sector circular, figuras compuestas, teorema de Pitágoras}}

% colors de la web (bloc d'Àlgebra: blau)
\definecolor{fmnavy}{HTML}{1B2748}
\definecolor{fmblue}{HTML}{3D6AA0}
\definecolor{fmbluel}{HTML}{4E81BD}
\definecolor{fmyellow}{HTML}{FFC94D}
\definecolor{fmgold}{HTML}{E9A919}
\definecolor{fmaccent}{HTML}{FF7A50}
\definecolor{fmdef}{HTML}{1D4F73}\definecolor{fmdefbg}{HTML}{EEF4F9}
\definecolor{fmex}{HTML}{A3203A}\definecolor{fmexbg}{HTML}{FBF0F2}
\definecolor{fmsf}{HTML}{C2410C}
\definecolor{fmok}{HTML}{15803D}
\definecolor{fmmuted}{HTML}{5D6685}
\definecolor{fmsurface}{HTML}{EEF1F8}
\definecolor{fmA}{HTML}{2F6FD6}\definecolor{fmB}{HTML}{EA6A10}\definecolor{fmR}{HTML}{E11D48}\definecolor{fmK}{HTML}{16A34A}\definecolor{fmV}{HTML}{7C3AED}

% lletres
\newcommand{\comic}{\fontfamily{ComicNeue-TLF}\selectfont}
\newcommand{\heavy}{\fontseries{eb}\selectfont}
\newcommand{\semi}{\fontseries{sb}\selectfont}
\DeclareUnicodeCharacter{2212}{-}

\setlength{\parindent}{0pt}
\setlength{\parskip}{7pt}
\setlength{\lineskip}{4pt}\setlength{\lineskiplimit}{4pt}
\setlength{\emergencystretch}{3em}
\setlist{itemsep=3pt,topsep=3pt,leftmargin=1.4em}
\setcounter{secnumdepth}{1}
\renewcommand{\thesection}{8.\arabic{section}}
\titleformat{\section}[hang]{\color{fmnavy}\fontsize{23}{27}\heavy\selectfont}{\tcbox[on line,colback=fmblue,colframe=fmblue,arc=3mm,boxsep=2pt,left=6pt,right=6pt,top=3pt,bottom=3pt]{\color{white}\thesection}}{0.45em}{}[\vspace{5pt}{\color{fmblue}\titlerule[2pt]}]
\titlespacing*{\section}{0pt}{4pt}{6pt}
\titleformat{\subsection}{\color{fmblue}\Large\semi}{}{0pt}{\faCaretRight\ }
\titlespacing*{\subsection}{0pt}{8pt}{2pt}

% frase de cada apartat, «a mà»
\newcommand{\frase}[1]{\par\vspace{-2pt}{\comic\Large\color{fmbluel}\faCommentDots[regular]\ \ #1\par}\vspace{2pt}}
\newcommand{\seccio}[3]{\Needspace{16\baselineskip}\section{#2}\hypertarget{#1}{}\frase{#3}}
% notes a mà dins dels exemples
\newcommand{\nota}[1]{{\comic\small\color{fmex!85!black}#1}}
\newcommand{\notab}[1]{{\comic\small\color{fmblue}#1}}
% bé / malament
\newcommand{\be}[1]{{\color{fmok}\faCheckCircle}\ \ensuremath{#1}}
\newcommand{\mal}[1]{{\color{fmR}\faTimesCircle}\ {\color{fmR}\ensuremath{#1}}}
% pas numerat
\newcommand{\pas}[1]{\tikz[baseline=(p.base)]\node[circle,fill=fmblue,text=white,inner sep=1.2pt,minimum size=1.25em,font=\small\heavy](p){#1};}
% etiqueta de color
\newcommand{\chip}[2]{\tcbox[on line,colback=#1!14,colframe=#1!60,arc=2mm,boxsep=1pt,left=5pt,right=5pt,top=1pt,bottom=1pt]{#2}}
\newcommand{\fmkey}[1]{\colorbox{fmyellow!55}{\strut #1}}
\newcommand{\cA}[1]{{\color{fmA}#1}}\newcommand{\cB}[1]{{\color{fmB}#1}}\newcommand{\cK}[1]{{\color{fmK}#1}}\newcommand{\cR}[1]{{\color{fmR}#1}}\newcommand{\cV}[1]{{\color{fmV}#1}}

% caixes
\tcbset{fmbase/.style={enhanced,arc=3mm,boxrule=1.2pt,left=4mm,right=4mm,top=2mm,bottom=2.5mm,before skip=8pt,after skip=8pt,fonttitle=\heavy\large,before upper=\raggedright}}
\newtcolorbox{idea}[1][Idea clave]{fmbase,colback=fmdefbg,colframe=fmdef,coltitle=white,title={\faKey\ \ #1}}
\newtcolorbox{exemple}[1][Ejemplo]{fmbase,colback=fmexbg,colframe=fmex,colbacktitle=fmexbg,coltitle=fmex,title={\faPencil*\ \ #1},titlerule=0.5pt,titlerule style={fmex!40,dashed}}
\newtcolorbox{compte}[1][¡Cuidado!]{fmbase,colback=fmsf!6,colframe=fmsf,coltitle=white,fonttitle=\comic\bfseries\Large,title={\faExclamationTriangle\ \ #1}}
\newtcolorbox{truc}[1][Truco]{fmbase,colback=fmyellow!20,colframe=fmgold,colbacktitle=fmyellow,coltitle=fmnavy,fonttitle=\comic\bfseries\Large,title={\faLightbulb[regular]\ \ #1}}

% dibuix de la targeta de la unitat, en blanc
\newcommand{\fmcard}{\begin{tikzpicture}[x=0.062cm,y=-0.062cm,white,line width=3pt,line cap=round,line join=round]
  \begin{scope}[transparency group,opacity=0.15]
  \node[anchor=base west,inner sep=0pt,font=\fontsize{28}{28}\selectfont] at (14,44) {$A=b\cdot h$};
  \draw (140,110) rectangle (230,50); \draw (185,35) circle (25);
  \end{scope}
\end{tikzpicture}}
\newcommand{\fmbrand}{\raisebox{-0.25\height}{\includegraphics[height=0.5cm]{logo-fentmates.pdf}}}

% regla de Ruffini: #1 = a, #2 = coeficients, #3 = fila del mig (sense el primer), #4 = fila de baix, #5 = nombre de coeficients
\newcommand{\ruffini}[5]{\begin{tikzpicture}[x=1.25cm,y=0.85cm,every node/.style={font=\large}]
  \foreach \v [count=\i] in {#2} \node (t\i) at (\i,0) {$\v$};
  \foreach \v [count=\i] in {#3} {\pgfmathtruncatemacro\j{\i+1}\node (m\j) at (\j,-1) {$\v$};}
  \foreach \v [count=\i] in {#4} \node (b\i) at (\i,-2.1) {$\v$};
  \node[text=fmB] at (0,-1) {$#1$};
  \draw[line width=1.1pt,fmnavy] (0.45,0.5)--(0.45,-1.5) (0.45,-1.5)--(#5+0.5,-1.5);
  \node[draw=fmR,line width=1.3pt,rounded corners=2mm,minimum size=0.75cm] at (#5,-2.1) {};
  \draw[-{Latex[length=2.2mm]},fmK,line width=1.1pt] (t1.south) -- (b1.north);
  \pgfmathtruncatemacro\last{#5-1}
  \foreach \i in {1,...,\last} {\pgfmathtruncatemacro\j{\i+1}
    \draw[-{Latex[length=2.2mm]},fmB,line width=1.1pt,shorten >=1pt] (b\i.north east) -- (m\j.south west);
    \draw[-{Latex[length=2.2mm]},fmA!70,line width=1pt,dashed,shorten >=2pt] ([xshift=8pt]t\j.south) -- ([xshift=8pt]b\j.north);}
\end{tikzpicture}}

% capçalera i peu
\pagestyle{fancy}\fancyhf{}
\fancyhead[L]{\small\fmbrand}
\fancyhead[R]{\small\color{fmmuted}1.º ESO · Unidad 8: Áreas y perímetros}
\fancyfoot[L]{\scriptsize\color{fmmuted}\href{https://creativecommons.org/licenses/by-nc-sa/4.0/deed.es}{\ccbyncsa}\ \ FentMates · fentmates.org · CC BY-NC-SA 4.0}
\fancyfoot[R]{\small\color{fmmuted}\thepage}
\renewcommand{\headrulewidth}{0.4pt}\renewcommand{\headrule}{\color{fmmuted!40}\hrule height \headrulewidth}
\fancypagestyle{primera}{\fancyhf{}\renewcommand{\headrulewidth}{0pt}
  \fancyfoot[L]{\scriptsize\color{fmmuted}\href{https://creativecommons.org/licenses/by-nc-sa/4.0/deed.es}{\ccbyncsa}\ \ FentMates · fentmates.org · CC BY-NC-SA 4.0}
  \fancyfoot[R]{\small\color{fmmuted}\thepage}}

% targeta de l'índex
\newcommand{\itile}[4]{\begin{tcolorbox}[enhanced,colback=white,colframe=fmblue!35,arc=3mm,boxrule=1pt,boxsep=0pt,left=6pt,right=6pt,top=5pt,bottom=5pt,before skip=0pt,after skip=7pt]
  \hyperlink{#1}{\begin{minipage}[c]{0.9cm}\centering{\color{fmblue}\LARGE #3}\end{minipage}\hspace{6pt}%
  \begin{minipage}[c]{0.78\linewidth}{\color{fmblue}\heavy #2}\\[-2pt]{\color{fmnavy}\semi\large #4}\end{minipage}}\end{tcolorbox}}

% taula amb vores arrodonides: el contingut es retalla amb la forma (sense @{} a les vores ni entre columnes)
\newsavebox\tabclipbox
\newcommand{\taulaclip}[3][3mm]{\sbox\tabclipbox{#3}\begin{tikzpicture}[baseline=(tcn.base)]
  \begin{scope}\clip[rounded corners=#1] (0,-\dp\tabclipbox) rectangle (\wd\tabclipbox,\ht\tabclipbox);
  \node[inner sep=0pt,outer sep=0pt,anchor=base west] (tcn) at (0,0) {\usebox\tabclipbox};\end{scope}
  \draw[#2,rounded corners=#1] (0,-\dp\tabclipbox) rectangle (\wd\tabclipbox,\ht\tabclipbox);
\end{tikzpicture}}
\newcommand{\taularod}[3][fmblue]{\taulaclip{#1,line width=1.2pt}{\renewcommand{\arraystretch}{1.5}\begin{tabular}{#2}#3\end{tabular}}}
% balança
\newcommand{\balanca}[3]{\begin{tikzpicture}[scale=#3,transform shape,every node/.style={anchor=center}]
  \fill[fmnavy] (-0.45,-1.9) rectangle (0.45,-1.75);
  \fill[fmnavy!80] (-0.12,-1.75) rectangle (0.12,0);
  \fill[fmgold] (0,0.05) circle (0.13);
  \draw[fmnavy,line width=2.4pt,line cap=round] (-2,0) -- (2,0);
  \foreach \s in {-1,1} {
    \draw[fmmuted,line width=0.8pt] (\s*2,0) -- (\s*2-0.75,-1.05) (\s*2,0) -- (\s*2+0.75,-1.05);
    \fill[fmyellow!70,draw=fmgold,line width=1pt] (\s*2-0.95,-1.05) arc[start angle=180,end angle=360,x radius=0.95,y radius=0.32] -- cycle;}
  \node[font=\Large,fill=fmA!15,draw=fmA,rounded corners=1.5mm,inner sep=3pt] at (-2,-0.68) {$#1$};
  \node[font=\Large,fill=fmB!15,draw=fmB,rounded corners=1.5mm,inner sep=3pt] at (2,-0.68) {$#2$};
\end{tikzpicture}}
% eixos amb quadrícula (dins d'un tikzpicture): xmin, xmax, ymin, ymax
\newcommand{\eixos}[4]{\draw[fmmuted!22,line width=0.4pt] (#1,#3) grid (#2,#4);
  \draw[-{Latex[length=2mm]},fmmuted,line width=0.8pt] (#1-0.2,0)--(#2+0.4,0) node[right,font=\small]{$x$};
  \draw[-{Latex[length=2mm]},fmmuted,line width=0.8pt] (0,#3-0.2)--(0,#4+0.4) node[above,font=\small]{$y$};}
% taula de traduccions (diccionari) amb les vores arrodonides
\newcommand{\dic}[2]{#1 & {\large$#2$}\\}
\newcommand{\dicratio}{0.62}
\newcommand{\dichead}{\begin{tabular}{@{}>{\raggedright}p{\dimexpr\dicratio\linewidth-12pt}>{\centering\arraybackslash}p{\dimexpr\linewidth-\dicratio\linewidth-12pt}@{}}\faComment[regular]\ Lo decimos así\dots & \dots\ así\end{tabular}}
\newenvironment{dicbox}{\begin{tcolorbox}[enhanced,arc=3mm,boxrule=1.2pt,colframe=fmblue,colback=white,coltitle=white,fonttitle=\heavy,
  title={\dichead},left=0pt,right=0pt,top=0pt,bottom=1.5mm,boxsep=0pt,lefttitle=6pt,righttitle=6pt,before skip=0pt,after skip=0pt]\rowcolors{1}{white}{fmsurface}%
  \begin{tabular}{@{\hspace{6pt}}>{\raggedright}p{\dimexpr\dicratio\linewidth-12pt}>{\centering\arraybackslash}p{\dimexpr\linewidth-\dicratio\linewidth-12pt}@{\hspace{6pt}}}}
  {\end{tabular}\end{tcolorbox}}
% resum
\newcommand{\rtile}[3]{\begin{tcolorbox}[enhanced,colback=white,colframe=fmblue!45,arc=3mm,boxrule=1pt,left=3mm,right=3mm,top=2.5mm,bottom=2.5mm,fontupper=\large,
  title={\tcbox[on line,colback=fmblue,colframe=fmblue,arc=2mm,boxsep=0pt,left=4pt,right=4pt,top=1pt,bottom=1pt]{\color{white}\small\heavy #1}\ \ #2},coltitle=fmnavy,colbacktitle=fmdefbg,fonttitle=\semi,toptitle=0.5mm,bottomtitle=0.5mm]#3\end{tcolorbox}}
% portada: #1 subtítol a mà
\newcommand{\portada}[1]{\thispagestyle{primera}
\begin{tikzpicture}[remember picture,overlay]
  \fill[fmblue] (current page.north west) rectangle ([yshift=-9.4cm]current page.north east);
  \begin{scope}
    \clip (current page.north west) rectangle ([yshift=-9.4cm]current page.north east);
    \node[anchor=north west,inner sep=0pt] at ([xshift=6.6cm,yshift=-1.3cm]current page.north west) {\fmcard};
  \end{scope}
  \node[anchor=north west,inner sep=0pt] at ([xshift=1.6cm,yshift=-1.1cm]current page.north west) {\tcbox[colback=white,colframe=white,arc=2.5mm,boxsep=0pt,left=7pt,right=7pt,top=5pt,bottom=5pt]{\includegraphics[height=0.75cm]{logo-fentmates.pdf}}};
\end{tikzpicture}
\vspace*{0.55cm}
\tcbox[colback=white!15!fmblue,colframe=white!15!fmblue,arc=2.5mm,boxsep=0pt,left=6pt,right=6pt,top=3pt,bottom=3pt]{\color{white}\small\bfseries Apuntes de teoría · 1.º ESO · Matemáticas}

\vspace{2pt}
{\raggedright\color{white}\fontsize{30}{36}\heavy\selectfont{\color{fmyellow}Unidad 8}\\[2pt]Áreas y perímetros\par}
\vspace{8pt}
{\color{white}\comic\Large #1\par}
\vspace*{\dimexpr 7.4cm-\pagetotal\relax}
}
\newcommand{\llegenda}{\par\vspace{16pt}
{\color{fmnavy}\heavy\Large\faPalette\ \ Cómo se leen estos apuntes}\par\smallskip
\begin{tcbraster}[raster columns=4,raster equal height,raster column skip=8pt,raster before skip=12pt]
\tcbox[colback=fmdefbg,colframe=fmdef,arc=3mm,boxrule=1.2pt,halign=center,width=\linewidth]{{\color{fmdef}\Large\faKey}\\[2pt]{\color{fmdef}\heavy Idea clave}\\[-1pt]{\small lo esencial}}
\tcbox[colback=fmexbg,colframe=fmex,arc=3mm,boxrule=1.2pt,halign=center,width=\linewidth]{{\color{fmex}\Large\faPencil*}\\[2pt]{\color{fmex}\heavy Ejemplo}\\[-1pt]{\small paso a paso}}
\tcbox[colback=fmsf!6,colframe=fmsf,arc=3mm,boxrule=1.2pt,halign=center,width=\linewidth]{{\color{fmsf}\Large\faExclamationTriangle}\\[2pt]{\color{fmsf}\heavy ¡Cuidado!}\\[-1pt]{\small errores típicos}}
\tcbox[colback=fmyellow!20,colframe=fmgold,arc=3mm,boxrule=1.2pt,halign=center,width=\linewidth]{{\color{fmgold!80!black}\Large\faLightbulb[regular]}\\[2pt]{\color{fmnavy}\heavy Truco}\\[-1pt]{\small para recordarlo}}
\end{tcbraster}}
\newcommand{\resumtitol}{\clearpage
\hypertarget{resum}{}
{\color{fmnavy}\fontsize{23}{27}\heavy\selectfont\faClipboardList\ \ En resumen}\par\vspace{5pt}{\color{fmblue}\titlerule[2pt]}
\frase{Toda la unidad en una página.}}
% els 5 passos per a resoldre problemes: #2 = text del pas 2
\newcommand{\passosproblema}[2]{\begin{center}
\begin{tikzpicture}[s/.style={draw=fmblue,fill=fmdefbg,rounded corners=3mm,line width=1.2pt,minimum width=2.75cm,minimum height=2.45cm,align=center,font=\small},
  ic/.style={font=\LARGE,text=fmblue}]
  \foreach \i/\ic/\t in {1/\faBookOpen/{\textbf{Leemos}\\y entendemos},2/\faQuestionCircle/{#1},3/\faEquals/{#2},4/\faCogs/{\textbf{Resolvemos}},5/\faCheckDouble/{\textbf{Comprobamos}\\y respondemos}} {
    \node[s] (s\i) at (3.3*\i,0) {};
    \node[ic] at ([yshift=-0.5cm]s\i.north) {\ic};
    \node[align=center,font=\small,anchor=north] at ([yshift=-0.95cm]s\i.north) {\t};
    \node[circle,fill=fmblue,text=white,font=\small\heavy,inner sep=1.5pt,minimum size=1.3em] at (s\i.north west) {\i};}
  \foreach \i [evaluate=\i as \j using int(\i+1)] in {1,...,4} \draw[-{Latex[length=3mm]},fmmuted,line width=1.5pt] (s\i) -- (s\j);
\end{tikzpicture}
\end{center}}

\newcommand{\lo}{\mathopen{]}}\newcommand{\ro}{\mathclose{[}}
% punts ple i buit, nombres dels eixos i taula de valors
\newcommand{\pple}[2][fmblue]{\fill[#1] #2 circle (3.6pt);}
\newcommand{\nums}[4]{\foreach \k in {#1,...,#2} {\ifnum\k=0\else\node[below,font=\scriptsize,text=fmmuted] at (\k,0) {$\k$};\fi}
  \foreach \k in {#3,...,#4} {\ifnum\k=0\else\node[left,font=\scriptsize,text=fmmuted] at (0,\k) {$\k$};\fi}}
\newcommand{\tv}[3][fmblue]{\taulaclip{#1,line width=1.2pt}{\renewcommand{\arraystretch}{1.45}%
  \begin{tabular}{>{\columncolor{#1}\color{white}\heavy}c*{12}{w{c}{0.8cm}}}#2\\\rowcolor{#1!8}\cellcolor{#1}#3\end{tabular}}}
% etiqueta d'una recta o corba
\newcommand{\etq}[3]{\node[fill=white,inner sep=1.5pt,text=#1,font=\small] at #2 {#3};}

% interval sobre la recta: #1 esquerre, #2 dret, #3 tipus esquerre (o/c/a), #4 tipus dret (o/c/a), #5 color
\def\intvdot#1#2#3{\ifx#2o\filldraw[fill=white,draw=#3,line width=1.3pt] (#1,0) circle (3.3pt);\fi\ifx#2c\fill[#3] (#1,0) circle (3.5pt);\fi}
\newcommand{\intv}[5]{\begin{tikzpicture}[x=0.4cm,baseline=-3pt]
  \draw[fmmuted,line width=0.8pt] (-6,0)--(7,0);
  \foreach \k in {-5,...,6} \draw[fmmuted,line width=0.6pt] (\k,-0.08)--(\k,0.08);
  \ifx#3a\def\xa{-6}\else\def\xa{#1}\fi
  \ifx#4a\def\xb{7}\else\def\xb{#2}\fi
  \ifx#4a\draw[#5,line width=2.6pt,-{Latex[length=2.6mm]}] (\xa,0)--(\xb,0);\else\ifx#3a\draw[#5,line width=2.6pt,{Latex[length=2.6mm]}-] (\xa,0)--(\xb,0);\else\draw[#5,line width=2.6pt] (\xa,0)--(\xb,0);\fi\fi
  \intvdot{#1}{#3}{#5}\intvdot{#2}{#4}{#5}
  \ifx#3a\else\node[below=2pt,font=\scriptsize] at (#1,0) {$#1$};\fi
  \ifx#4a\else\node[below=2pt,font=\scriptsize] at (#2,0) {$#2$};\fi
\end{tikzpicture}}

\newcommand{\angrect}[3]{\draw[fmmuted,line width=0.8pt] ($#2!0.22cm!#1$) -- ($#2!0.22cm!#1!0.22cm!90:#1$) -- ($#2!0.22cm!#3$);}
\newcommand{\figura}[3]{\begin{tcolorbox}[fmbase,colback=fmdefbg,colframe=fmdef,coltitle=white,title={#1},before upper=\centering,left=2mm,right=2mm]
#2\par\smallskip #3
\end{tcolorbox}}


\begin{document}
\portada{Medimos contornos y superficies de las figuras planas.}

{\color{fmnavy}\heavy\Large\faMapSigns\ \ ¿Qué aprenderás?}\par\smallskip
\begin{minipage}[t]{0.485\linewidth}
\itile{sec1}{8.1}{\faRulerCombined}{Perímetro, área y unidades}
\itile{sec2}{8.2}{\faRulerHorizontal}{Perímetro de los polígonos}
\itile{sec3}{8.3}{\faSquare}{Cuadrado y rectángulo}
\itile{sec4}{8.4}{\faIcon{play}}{Romboide y triángulo}
\itile{sec5}{8.5}{\faIcon{gem}}{Rombo y trapecio}
\itile{sec6}{8.6}{\faIcon{stop-circle}}{Polígonos regulares}
\end{minipage}\hfill
\begin{minipage}[t]{0.485\linewidth}
\itile{sec7}{8.7}{\faIcon{circle-notch}}{La circunferencia}
\itile{sec8}{8.8}{\faChartPie}{Área del círculo}
\itile{sec9}{8.9}{\faPuzzlePiece}{Figuras compuestas}
\itile{sec10}{8.10}{\faIcon{square}}{El teorema de Pitágoras}
\itile{sec11}{8.11}{\faSearch}{Pitágoras en áreas y perímetros}
\itile{resum}{\phantom{8.11}}{\faClipboardList}{En resumen}
\end{minipage}
\llegenda

% ============================================================ 8.1
\clearpage
\seccio{sec1}{Perímetro y área. Unidades de medida}{El contorno y la superficie.}

\vfill
\begin{minipage}[c]{0.55\linewidth}
\begin{idea}[Perímetro y área]
\textbf{Perímetro}: la longitud del \textbf{contorno}, el camino de una vuelta entera. En m, cm, km\dots\\[3pt]
\textbf{Área}: la medida de la \textbf{superficie}, cuánto sitio hay dentro. En unidades cuadradas: 1 cm$^2$ es un cuadrado de 1 cm de lado.
\end{idea}
\end{minipage}\hfill
\begin{minipage}[c]{0.42\linewidth}\centering
\begin{tikzpicture}[scale=0.5]
  \foreach \p in {(0,0),(1,0),(2,0),(3,0),(0,1),(1,1),(1,2)} \filldraw[fill=fmA!35,draw=fmA] \p rectangle ++(1,1);
  \draw[fmR,line width=1.6pt] (0,0) -- (4,0) -- (4,1) -- (2,1) -- (2,3) -- (1,3) -- (1,2) -- (0,2) -- cycle;
  \draw[fmmuted!40] (-0.5,-0.5) grid (4.5,3.5);
\end{tikzpicture}\\
{\small$A=7$ cm$^2$ \ (cuadraditos) \qquad \cR{$P=14$ cm} \ (contorno)}
\end{minipage}

\vfill
\begin{compte}[Área y perímetro son cosas distintas]
Un rectángulo de $6\times2$ y uno de $4\times3$ tienen la misma área (12 cm$^2$) pero perímetros distintos: 16 cm y 14 cm.
\end{compte}

\vfill
\begin{center}\small
\taularod{>{\heavy}l*{7}{>{\hspace{4pt}}c<{\hspace{4pt}}}}{
\rowcolor{fmblue}\color{white}Longitud & \color{white}km & \color{white}hm & \color{white}dam & \color{white}\textbf{m} & \color{white}dm & \color{white}cm & \color{white}mm\\
 & 1\,000 m & 100 m & 10 m & 1 m & 0,1 m & 0,01 m & 0,001 m
}\\[4pt]
{\comic\color{fmmuted}de 10 en 10: \ hacia la derecha $\times10$, \ hacia la izquierda $:10$}\\[8pt]
\taularod{>{\heavy}l*{7}{>{\hspace{4pt}}c<{\hspace{4pt}}}}{
\rowcolor{fmK!80!black}\color{white}Superficie & \color{white}km$^2$ & \color{white}hm$^2$ & \color{white}dam$^2$ & \color{white}\textbf{m}$^2$ & \color{white}dm$^2$ & \color{white}cm$^2$ & \color{white}mm$^2$\\
 & & hectárea (ha) & área (a) & 1 m$^2$ & 0,01 m$^2$ & 0,0001 m$^2$ &
}\\[4pt]
{\comic\color{fmmuted}de 100 en 100: \ hacia la derecha $\times100$, \ hacia la izquierda $:100$}
\end{center}

\vfill
\begin{minipage}[c]{0.3\linewidth}\centering
\begin{tikzpicture}[scale=0.24]
  \draw[fmmuted!50] (0,0) grid (10,10); \draw[fmK,line width=1.4pt] (0,0) rectangle (10,10);
  \node[below,font=\scriptsize] at (5,0) {1 m $=10$ dm};
\end{tikzpicture}
\end{minipage}\hfill
\begin{minipage}[c]{0.67\linewidth}
\begin{truc}[¿Por qué de 100 en 100?]
Un cuadrado de 1 m de lado mide 10 dm de lado: caben $10\times10=\boxed{100}$ cuadraditos de 1 dm$^2$. \ $1\text{ m}^2=100\text{ dm}^2$.\\[2pt]
Campos: $1\text{ ha}=1\text{ hm}^2=10\,000\text{ m}^2$ \ y \ $1\text{ a}=100\text{ m}^2$.
\end{truc}
\end{minipage}

\vfill
\begin{exemple}[Ejemplos]
\begin{minipage}[t]{0.48\linewidth}
$3{,}5\text{ km}=3\,500\text{ m}$ \qquad $245\text{ cm}=2{,}45\text{ m}$\\[4pt]
$4\text{ m}=400\text{ cm}$ \qquad $80\text{ mm}=8\text{ cm}$
\end{minipage}\hfill
\begin{minipage}[t]{0.49\linewidth}
$3\text{ m}^2=300\text{ dm}^2=30\,000\text{ cm}^2$\\[4pt]
$3{,}5\text{ ha}=35\,000\text{ m}^2$ \qquad $500\text{ cm}^2=5\text{ dm}^2$
\end{minipage}
\end{exemple}

\vfill
% ============================================================ 8.2
\clearpage
\seccio{sec2}{Perímetro de los polígonos}{Sumamos todos los lados.}

\vfill
\begin{idea}[Perímetro de un polígono]
La suma de las longitudes de \textbf{todos} los lados. Si hay iguales, multiplicamos.
\end{idea}

\vfill
\begin{tcbraster}[raster columns=3,raster equal height=rows,raster column skip=8pt,raster row skip=8pt]
\figura{Triángulo}{\begin{tikzpicture}[scale=0.4]\filldraw[fill=fmK!20,draw=fmK,line width=1.3pt] (0,0) -- (4.4,0) -- (1.4,2.6) -- cycle;\node[below,font=\scriptsize] at (2.2,0) {$a$};\node[left,font=\scriptsize] at (0.7,1.3) {$b$};\node[right,font=\scriptsize] at (2.9,1.3) {$c$};\end{tikzpicture}}{\large$P=a+b+c$}
\figura{Cuadrado}{\begin{tikzpicture}[scale=0.4]\filldraw[fill=fmA!20,draw=fmA,line width=1.3pt] (0,0) rectangle (2.6,2.6);\node[below,font=\scriptsize] at (1.3,0) {$l$};\end{tikzpicture}}{\large$P=4\cdot l$}
\figura{Rectángulo}{\begin{tikzpicture}[scale=0.4]\filldraw[fill=fmA!20,draw=fmA,line width=1.3pt] (0,0) rectangle (4.2,2.6);\node[below,font=\scriptsize] at (2.1,0) {$b$};\node[right,font=\scriptsize] at (4.2,1.3) {$h$};\end{tikzpicture}}{\large$P=2\cdot b+2\cdot h$}
\figura{Romboide}{\begin{tikzpicture}[scale=0.4]\filldraw[fill=fmB!20,draw=fmB,line width=1.3pt] (0,0) -- (3.6,0) -- (4.8,2.6) -- (1.2,2.6) -- cycle;\node[below,font=\scriptsize] at (1.8,0) {$b$};\node[left,font=\scriptsize] at (0.5,1.4) {$c$};\end{tikzpicture}}{\large$P=2\cdot b+2\cdot c$}
\figura{Rombo}{\begin{tikzpicture}[scale=0.4]\filldraw[fill=fmB!20,draw=fmB,line width=1.3pt] (0,1.3) -- (2.2,0) -- (4.4,1.3) -- (2.2,2.6) -- cycle;\node[below left,font=\scriptsize] at (1.1,0.65) {$l$};\end{tikzpicture}}{\large$P=4\cdot l$}
\figura{Polígono regular}{\begin{tikzpicture}[scale=0.55]\filldraw[fill=fmV!20,draw=fmV,line width=1.3pt] (90:1.2) -- (162:1.2) -- (234:1.2) -- (306:1.2) -- (18:1.2) -- cycle;\node[below,font=\scriptsize] at (270:1) {$l$};\end{tikzpicture}}{\large$P=n\cdot l$ \ \small($n$ lados)}
\end{tcbraster}

\vfill
\begin{compte}[Todos los lados en la misma unidad]
Un rectangle d'1,2 m $\times$ 80 cm: \ $P=2\cdot120+2\cdot80=400$ cm $=4$ m. \ \mal{2\cdot1{,}2+2\cdot80}
\end{compte}

\vfill
\begin{exemple}[Ejemplos]
\begin{minipage}[t]{0.48\linewidth}
Triángulo de lados 5, 7 y 9 cm: \ $P=\boxed{21\text{ cm}}$\\[4pt]
Cuadrado de 6,5 m: \ $P=4\cdot6{,}5=\boxed{26\text{ m}}$\\[4pt]
Rectángulo de 12 m $\times$ 7 m: \ $P=24+14=\boxed{38\text{ m}}$
\end{minipage}\hfill
\begin{minipage}[t]{0.49\linewidth}
Hexágono regular de 8 cm: \ $P=6\cdot8=\boxed{48\text{ cm}}$\\[4pt]
Pentágono regular de $P=45$ cm: \ $l=45:5=\boxed{9\text{ cm}}$\\[4pt]
Rombo de $P=34$ cm: \ $l=34:4=\boxed{8{,}5\text{ cm}}$
\end{minipage}
\end{exemple}

\vfill
\begin{exemple}[Ejemplo: \ queremos cercar un huerto rectangular de 15 m $\times$ 8 m con tres vueltas de alambre. ¿Cuánto alambre hace falta?]
Perímetro: \ $2\cdot15+2\cdot8=46$ m \qquad Tres vueltas: \ $3\cdot46=\boxed{138\text{ m}}$ de alambre
\end{exemple}

\vfill
% ============================================================ 8.3
\clearpage
\seccio{sec3}{Área del cuadrado y del rectángulo}{Contamos cuántas filas de cuadraditos caben.}

\vfill
\begin{minipage}[c]{0.45\linewidth}\centering
\begin{tikzpicture}[scale=0.55]
  \foreach \i in {0,...,4} \foreach \j in {0,1,2} \filldraw[fill=fmA!30,draw=fmA] (\i,\j) rectangle ++(1,1);
  \draw[fmA,line width=1.5pt] (0,0) rectangle (5,3);
  \node[below,font=\small] at (2.5,0) {$b=5$ cm}; \node[right,font=\small] at (5,1.5) {$h=3$ cm};
\end{tikzpicture}\\[3pt]
{\comic\small\color{fmmuted}3 filas de 5 cuadraditos: $5\cdot3=15$ cm$^2$}
\end{minipage}\hfill
\begin{minipage}[c]{0.52\linewidth}
\begin{tcbraster}[raster columns=2,raster equal height,raster column skip=8pt]
\begin{tcolorbox}[fmbase,colback=fmA!8,colframe=fmA,coltitle=white,title={Rectángulo},before upper=\centering]
{\Large$A=b\cdot h$}
\end{tcolorbox}
\begin{tcolorbox}[fmbase,colback=fmK!8,colframe=fmK,coltitle=white,title={Cuadrado},before upper=\centering]
{\Large$A=l^2$}
\end{tcolorbox}
\end{tcbraster}
\small También con decimales: $4{,}5\text{ m}\times3\text{ m}$, \ $A=13{,}5$ m$^2$.
\end{minipage}

\vfill
\begin{tcbraster}[raster columns=2,raster equal height,raster column skip=10pt]
\begin{tcolorbox}[fmbase,colback=fmdefbg,colframe=fmdef,coltitle=white,title={Conocemos el área y la base}]
La altura: \ $h=A:b$\par\smallskip
Rectángulo de 60 cm$^2$ y base 12 cm: \ $h=60:12=5$ cm
\end{tcolorbox}
\begin{tcolorbox}[fmbase,colback=fmdefbg,colframe=fmdef,coltitle=white,title={Conocemos el área del cuadrado}]
El lado: \ $l=\sqrt A$\par\smallskip
Cuadrado de 49 cm$^2$: \ $l=\sqrt{49}=7$ cm
\end{tcolorbox}
\end{tcbraster}

\vfill
\begin{compte}[La misma unidad, y el resultado al cuadrado]
Las dos medidas en la misma unidad: \ cm $\cdot$ cm $=$ cm$^2$. \ Un rectángulo de 2 m $\times$ 50 cm: \ $200\cdot50=10\,000$ cm$^2=1$ m$^2$.
\end{compte}

\vfill
\begin{exemple}[Ejemplos]
\begin{minipage}[t]{0.48\linewidth}
Rectángulo de 8 cm $\times$ 5 cm: \ $A=\boxed{40\text{ cm}^2}$\\[4pt]
Cuadrado de 9 m: \ $A=9^2=\boxed{81\text{ m}^2}$\\[4pt]
Campo de 120 m $\times$ 85 m: \ $A=10\,200$ m$^2\approx1$ ha
\end{minipage}\hfill
\begin{minipage}[t]{0.49\linewidth}
Cuadrado de perímetro 36 cm: \ $l=9$, \ $A=\boxed{81\text{ cm}^2}$\\[4pt]
Rectángulo de 54 m$^2$ y 9 m de largo: \ ancho $54:9=\boxed{6\text{ m}}$\\[4pt]
Cuadrado de 100 cm$^2$: \ $l=\boxed{10\text{ cm}}$
\end{minipage}
\end{exemple}

\vfill
\begin{exemple}[Ejemplo: \ queremos poner parqué en una habitación de 4 m $\times$ 3,5 m. El m$^2$ cuesta 25 €]
$A=4\cdot3{,}5=14$ m$^2$ \qquad Precio: \ $14\cdot25=\boxed{350\text{ €}}$
\end{exemple}

\vfill
% ============================================================ 8.4
\clearpage
\seccio{sec4}{Área del romboide y del triángulo}{Recortamos y pegamos: todo sale del rectángulo.}

\vfill
\begin{minipage}[c]{0.48\linewidth}\centering
\begin{tikzpicture}[scale=0.55]
  \filldraw[fill=fmB!25,draw=fmB,line width=1.4pt] (0,0) -- (5,0) -- (6.5,3) -- (1.5,3) -- cycle;
  \fill[fmR!30] (0,0) -- (1.5,0) -- (1.5,3) -- cycle; \draw[fmR,dashed,line width=1pt] (1.5,0) -- (1.5,3);
  \draw[fmR!70,dashed] (5,0) -- (6.5,0) -- (6.5,3);
  \draw[-{Latex[length=2mm]},fmR,line width=1pt] (0.8,1) to[bend right=25] (5.7,0.8);
  \node[below,font=\small] at (2.5,0) {$b$}; \node[left,font=\small,text=fmR] at (1.5,1.6) {$h$};
\end{tikzpicture}\\[2pt]
{\comic\small\color{fmmuted}el triángulo pasa al otro lado: un rectángulo}
\end{minipage}\hfill
\begin{minipage}[c]{0.49\linewidth}
\begin{tcolorbox}[fmbase,colback=fmB!8,colframe=fmB,coltitle=white,title={Romboide},before upper=\centering]
{\Large$A=b\cdot h$}\par\smallskip
\small la misma base y altura que el rectángulo
\end{tcolorbox}
\end{minipage}

\vfill
\begin{compte}[La altura no es el lado inclinado]
La altura es la distancia entre las dos bases, medida \textbf{perpendicularmente}.
\end{compte}

\vfill
\begin{minipage}[c]{0.48\linewidth}\centering
\begin{tikzpicture}[scale=0.55]
  \filldraw[fill=fmK!30,draw=fmK!80!black,line width=1.4pt] (0,0) -- (5,0) -- (1.5,3) -- cycle;
  \draw[fmK!80!black,dashed,line width=1pt] (5,0) -- (6.5,3) -- (1.5,3);
  \draw[fmR,dashed,line width=1pt] (1.5,0) -- (1.5,3); \node[left,font=\small,text=fmR] at (1.5,1.4) {$h$}; \node[below,font=\small] at (2.5,0) {$b$};
\end{tikzpicture}\\[2pt]
{\comic\small\color{fmmuted}la mitad de un romboide}
\end{minipage}\hfill
\begin{minipage}[c]{0.49\linewidth}
\begin{tcolorbox}[fmbase,colback=fmK!8,colframe=fmK,coltitle=white,title={Triángulo},before upper=\centering]
{\Large$A=\dfrac{b\cdot h}{2}$}\par\smallskip
\small una diagonal parte el romboide en dos triángulos iguales
\end{tcolorbox}
\end{minipage}

\vfill
\begin{tcbraster}[raster columns=3,raster equal height,raster column skip=8pt]
\figura{Acutángulo: dentro}{\begin{tikzpicture}[scale=0.38]\filldraw[fill=fmK!20,draw=fmK,line width=1.2pt] (0,0) -- (4,0) -- (1.5,2.8) -- cycle;\draw[fmR,dashed] (1.5,0) -- (1.5,2.8);\end{tikzpicture}}{}
\figura{Rectángulo: un cateto}{\begin{tikzpicture}[scale=0.38]\filldraw[fill=fmK!20,draw=fmK,line width=1.2pt] (0,0) -- (4,0) -- (0,2.8) -- cycle;\draw[fmR,line width=1.6pt] (0,0) -- (0,2.8);\end{tikzpicture}}{}
\figura{Obtusángulo: fuera}{\begin{tikzpicture}[scale=0.38]\filldraw[fill=fmK!20,draw=fmK,line width=1.2pt] (1.5,0) -- (4.5,0) -- (0,2.8) -- cycle;\draw[fmR,dashed] (0,0) -- (0,2.8); \draw[fmmuted,dashed] (0,0) -- (1.5,0);\end{tikzpicture}}{}
\end{tcbraster}

\vfill
\begin{exemple}[Ejemplos]
\begin{minipage}[t]{0.48\linewidth}
Romboide de $b=9$ cm y $h=4$ cm:\\ $A=\boxed{36\text{ cm}^2}$\\[4pt]
Triángulo de $b=10$ m y $h=6$ m:\\ $A=\dfrac{60}{2}=\boxed{30\text{ m}^2}$
\end{minipage}\hfill
\begin{minipage}[t]{0.49\linewidth}
Triángulo rectángulo de catetos 6 y 8 cm:\\ $A=\dfrac{6\cdot8}{2}=\boxed{24\text{ cm}^2}$\\[4pt]
Triángulo de 21 cm$^2$ y base 7 cm:\\ $h=\dfrac{2\cdot21}{7}=\boxed{6\text{ cm}}$
\end{minipage}
\end{exemple}

\vfill
% ============================================================ 8.5
\clearpage
\seccio{sec5}{Área del rombo y del trapecio}{Uno es la mitad de un rectángulo; el otro, la mitad de un romboide.}

\vfill
\begin{minipage}[c]{0.42\linewidth}\centering
\begin{tikzpicture}[scale=0.55]
  \draw[fmmuted,dashed,line width=1pt] (0,0) rectangle (5,3.4);
  \filldraw[fill=fmB!30,draw=fmB,line width=1.4pt] (0,1.7) -- (2.5,0) -- (5,1.7) -- (2.5,3.4) -- cycle;
  \draw[fmR,line width=1.1pt] (0,1.7) -- (5,1.7) node[pos=0.78,above,font=\small] {$D$}; \draw[fmA,line width=1.1pt] (2.5,0) -- (2.5,3.4) node[pos=0.25,left,font=\small] {$d$};
\end{tikzpicture}\\[2pt]
{\comic\small\color{fmmuted}la mitad del rectángulo $D\times d$}
\end{minipage}\hfill
\begin{minipage}[c]{0.55\linewidth}
\begin{tcolorbox}[fmbase,colback=fmB!8,colframe=fmB,coltitle=white,title={Rombo},before upper=\centering]
{\Large$A=\dfrac{D\cdot d}{2}$}\par\smallskip
\small $D$: diagonal mayor; \ $d$: menor. \ Como es un romboide, también $A=b\cdot h$.
\end{tcolorbox}
\end{minipage}

\vfill
\begin{minipage}[c]{0.42\linewidth}\centering
\begin{tikzpicture}[scale=0.42]
  \filldraw[fill=fmK!30,draw=fmK!80!black,line width=1.3pt] (0,0) -- (6,0) -- (4.5,3) -- (1,3) -- cycle;
  \filldraw[fill=fmK!12,draw=fmK!80!black,line width=1.1pt,dashed] (6,0) -- (4.5,3) -- (10.5,3) -- (9.5,0) -- cycle;
  \node[below,font=\small] at (3,0) {$B$}; \node[above,font=\small] at (2.7,3) {$b$}; \node[below,font=\small] at (7.75,0) {$b$}; \node[above,font=\small] at (7.5,3) {$B$};
  \draw[fmR,dashed] (1,0) -- (1,3); \node[left,font=\small,text=fmR] at (1,1.5) {$h$};
\end{tikzpicture}\\[2pt]
{\comic\small\color{fmmuted}dos trapecios: un romboide de base $B+b$}
\end{minipage}\hfill
\begin{minipage}[c]{0.55\linewidth}
\begin{tcolorbox}[fmbase,colback=fmK!8,colframe=fmK,coltitle=white,title={Trapecio},before upper=\centering]
{\Large$A=\dfrac{(B+b)\cdot h}{2}$}\par\smallskip
\small primero sumamos las bases (el paréntesis)
\end{tcolorbox}
\end{minipage}

\vfill
\begin{exemple}[Ejemplo: \ trapecio de $B=10$ cm, $b=6$ cm y $h=4$ cm]
$A=\dfrac{(10+6)\cdot4}{2}=\dfrac{16\cdot4}{2}=\dfrac{64}{2}=\boxed{32\text{ cm}^2}$
\end{exemple}

\vfill
\begin{exemple}[Ejemplos]
\begin{minipage}[t]{0.48\linewidth}
Rombo de diagonales 12 y 8 cm:\\ $A=\dfrac{96}{2}=\boxed{48\text{ cm}^2}$\\[4pt]
Rombo de lado 5 cm y altura 4,8 cm:\\ $A=5\cdot4{,}8=\boxed{24\text{ cm}^2}$
\end{minipage}\hfill
\begin{minipage}[t]{0.49\linewidth}
Trapecio de bases 9 y 5 m y altura 3 m:\\ $A=\dfrac{14\cdot3}{2}=\boxed{21\text{ m}^2}$\\[4pt]
Una ventana en forma de rombo de 80 cm y 50 cm:\\ $A=\boxed{2\,000\text{ cm}^2}$
\end{minipage}
\end{exemple}

\vfill
\begin{compte}[Los paréntesis del trapecio]
\be{\frac{(10+6)\cdot4}{2}=32} \qquad \mal{\frac{10+6\cdot4}{2}=17}: sin paréntesis multiplicaríamos solo la base pequeña.
\end{compte}

\vfill
% ============================================================ 8.6
\clearpage
\seccio{sec6}{Área de los polígonos regulares}{Todos los lados y todos los ángulos iguales.}

\vfill
\begin{minipage}[c]{0.42\linewidth}\centering
\begin{tikzpicture}[scale=1.25]
  \filldraw[fill=fmA!10,draw=fmnavy,line width=1.5pt] (0:2) -- (60:2) -- (120:2) -- (180:2) -- (240:2) -- (300:2) -- cycle;
  \fill[fmK!25] (0,0) -- (300:2) -- (330:1.732) -- cycle;
  \draw[fmB,line width=1.6pt] (0,0) -- (300:2) node[midway,left=1pt,font=\small\heavy] {$r$};
  \draw[fmK,line width=1.6pt] (0,0) -- (330:1.732) node[midway,above=1pt,font=\small\heavy] {$a$};
  \draw[fmR,line width=1.6pt] (330:1.732) -- (300:2) node[midway,right=1pt,font=\small\heavy] {$\frac l2$};
  \fill[fmnavy] (0,0) circle (1.6pt) node[left,font=\small] {$O$};
\end{tikzpicture}\\[2pt]
{\comic\small\color{fmmuted}hexágono regular}
\end{minipage}\hfill
\begin{minipage}[c]{0.55\linewidth}
\begin{idea}[Elementos]
\textbf{Centro $O$}: a la misma distancia de todos los vértices.\\[3pt]
\cB{\textbf{Radio $r$}}: del centro a un vértice.\\[3pt]
\cK{\textbf{Apotema $a$}}: del centro al punto medio de un lado (perpendicular).\\[4pt]
Radio, apotema y medio lado forman un \textbf{triángulo rectángulo}:\\[3pt]
{\centering\large$r^2=a^2+\left(\dfrac l2\right)^2$\par}
\end{idea}
\end{minipage}

\vfill
\begin{tcolorbox}[enhanced,colback=fmyellow!18,colframe=fmgold,arc=4mm,boxrule=1.6pt,before upper=\centering,top=2mm,bottom=2mm]
{\Large$A=\dfrac{P\cdot a}{2}$}\qquad perímetro por apotema, entre 2\par\smallskip
{\small un polígono de $n$ lados son $n$ triángulos de base $l$ y altura $a$}
\end{tcolorbox}

\vfill
\begin{minipage}[t]{0.5\linewidth}
\begin{exemple}[Ejemplo: \ pentágono de lado 8 cm y apotema 5,5 cm]
$P=5\cdot8=\boxed{40\text{ cm}}$\\[3pt]
$A=\dfrac{40\cdot5{,}5}{2}=\boxed{110\text{ cm}^2}$
\end{exemple}
\end{minipage}\hfill
\begin{minipage}[t]{0.47\linewidth}
\begin{truc}[El hexágono es especial]
Sus seis triángulos son \textbf{equiláteros}: el radio es igual al lado, \ $\boxed{r=l}$.
\end{truc}
\end{minipage}

\vfill
\begin{exemple}[Ejemplo: \ hexágono regular de 10 cm de lado]
\begin{minipage}[c]{0.6\linewidth}
$r=l=10$, \ medio lado 5: \ $a=\sqrt{10^2-5^2}=\sqrt{75}\approx8{,}66$ cm\\[4pt]
$A=\dfrac{60\cdot8{,}66}{2}\approx\boxed{259{,}8\text{ cm}^2}$
\end{minipage}\hfill
\begin{minipage}[c]{0.36\linewidth}\centering
\begin{tikzpicture}[scale=0.6]
  \draw[fmex,line width=1.3pt] (0:1.6) -- (60:1.6) -- (120:1.6) -- (180:1.6) -- (240:1.6) -- (300:1.6) -- cycle;
  \foreach \k in {0,60,...,300} \draw[fmex!50] (0,0) -- (\k:1.6);
\end{tikzpicture}
\end{minipage}
\end{exemple}

\vfill
\begin{exemple}[Ejemplo: \ cuadrado inscrito en una circunferencia de 5 cm de radio]
La diagonal es un diámetro, 10 cm: \ $l^2+l^2=100\ \to\ l=\sqrt{50}\approx\boxed{7{,}07\text{ cm}}$, \ $A=l^2=\boxed{50\text{ cm}^2}$
\end{exemple}

\vfill

\vfill
% ============================================================ 8.7
\clearpage
\seccio{sec7}{La circunferencia: longitud y arcos}{Vuelta entera, $2\pi r$; un arco, la parte que toca.}

\vfill
\begin{minipage}[c]{0.4\linewidth}\centering
\begin{tikzpicture}[scale=1.15]
  \fill[fmA!8] (0,0) circle (1.6);
  \draw[fmnavy,line width=1.6pt] (0,0) circle (1.6);
  \draw[fmB,line width=1.4pt] (0,0) -- (40:1.6) node[midway,above left=-1pt,font=\small\heavy] {$r$};
  \draw[fmA,line width=1.4pt] (-1.6,0) -- (1.6,0); \node[below,font=\small\heavy,text=fmA] at (-0.8,0) {$d$};
  \draw[fmK,line width=1.4pt] (200:1.6) -- (290:1.6); \node[font=\small,text=fmK!80!black] at (230:0.75) {cuerda};
  \draw[fmR,line width=3pt] (95:1.6) arc (95:150:1.6); \node[font=\small\heavy,text=fmR] at (122:1.95) {arc};
  \fill[fmnavy] (0,0) circle (1.4pt);
\end{tikzpicture}
\end{minipage}\hfill
\begin{minipage}[c]{0.57\linewidth}
\begin{idea}[Circunferencia y círculo]
La \textbf{circunferencia} es la línea de puntos a distancia $r$ (el \textbf{radio}) del centro. El \textbf{círculo} es la superficie que encierra.\\[3pt]
La \textbf{cuerda} une dos puntos; la más larga es el \textbf{diámetro}, $d=2r$. Un \textbf{arco} es un trozo de circunferencia.
\end{idea}
\end{minipage}

\vfill
\begin{tcbraster}[raster columns=2,raster equal height,raster column skip=10pt]
\begin{tcolorbox}[enhanced,colback=fmyellow!18,colframe=fmgold,arc=4mm,boxrule=1.6pt,before upper=\centering]
{\Large$L=2\pi r=\pi d$}\par\smallskip
{\small $\pi=3{,}14159\dots$: longitud entre diámetro, siempre el mismo}
\end{tcolorbox}
\begin{tcolorbox}[enhanced,colback=fmyellow!18,colframe=fmgold,arc=4mm,boxrule=1.6pt,before upper=\centering]
{\Large$L_{\text{arc}}=\dfrac{2\pi r\cdot n}{360}$}\par\smallskip
{\small un arco de $n^\circ$ es $\frac{n}{360}$ de la vuelta}
\end{tcolorbox}
\end{tcbraster}

\vfill
\begin{exemple}[Ejemplos]
Circunferencia de $r=4{,}5$ cm: \ $L=2\pi\cdot4{,}5=9\pi\approx\boxed{28{,}27\text{ cm}}$\\[5pt]
Arco de $120^\circ$ con $r=9$ cm: \ $L_{\text{arco}}=\dfrac{2\pi\cdot9\cdot120}{360}=6\pi\approx\boxed{18{,}85\text{ cm}}$\\[5pt]
Semicírculo de $r=5$ cm (arco $+$ diámetro): \ $P=5\pi+10\approx\boxed{25{,}71\text{ cm}}$\\[5pt]
Arco de 15,71 cm con $r=6$ cm: \ $n=\dfrac{15{,}71\cdot360}{12\pi}\approx\boxed{150^\circ}$
\end{exemple}

\vfill
\begin{exemple}[Ejemplo: \ la rueda de una bici mide 70 cm de diámetro. ¿Cuántas vueltas completas da en 1 km?]
\begin{minipage}[c]{0.2\linewidth}\centering
\begin{tikzpicture}[scale=0.55]\draw[fmex,line width=2pt] (0,0) circle (1.2);\foreach \k in {0,45,...,315} \draw[fmex!50] (0,0) -- (\k:1.2);\fill[fmex] (0,0) circle (2pt);\end{tikzpicture}
\end{minipage}\hfill
\begin{minipage}[c]{0.77\linewidth}
Una vuelta: \ $L=\pi\cdot70\approx219{,}91$ cm\\[3pt]
$1\text{ km}=100\,000$ cm: \ $\dfrac{100\,000}{219{,}91}\approx454{,}7$ \ $\to$ \ \boxed{454\text{ vueltas completas}}
\end{minipage}
\end{exemple}

\vfill
\begin{truc}[Calculadora]
Usa la tecla $\pi$ y redondea solo al final. Con $\pi\approx3{,}14$ el resultado cambia un poco.
\end{truc}

\vfill

\vfill
% ============================================================ 8.8
\clearpage
\seccio{sec8}{Área del círculo y figuras circulares}{Una pizza cortada en trozos se convierte en un rectángulo.}

\vfill
\begin{minipage}[c]{0.5\linewidth}\centering
\begin{tikzpicture}[scale=0.62]
  \foreach \i in {0,...,15} {
    \pgfmathsetmacro\xx{\i*0.39}
    \ifodd\i \filldraw[fill=fmB!30,draw=fmB,line width=0.8pt] (\xx-0.39,2.2) to[bend left=12] (\xx+0.39,2.2) -- (\xx,0) -- cycle;
    \else    \filldraw[fill=fmA!30,draw=fmA,line width=0.8pt] (\xx-0.39,0) to[bend right=12] (\xx+0.39,0) -- (\xx,2.2) -- cycle; \fi}
  \draw[decorate,decoration={brace,amplitude=4pt,mirror},line width=0.9pt] (-0.39,-0.25) -- (6.24,-0.25) node[midway,below=5pt,font=\small] {$\pi r$ (media circunferencia)};
  \draw[decorate,decoration={brace,amplitude=4pt},line width=0.9pt] (-0.55,0) -- (-0.55,2.2) node[midway,left=5pt,font=\small] {$r$};
\end{tikzpicture}
\end{minipage}\hfill
\begin{minipage}[c]{0.47\linewidth}
\begin{tcolorbox}[enhanced,colback=fmyellow!18,colframe=fmgold,arc=4mm,boxrule=1.6pt,before upper=\centering,top=3mm,bottom=3mm]
{\fontsize{24}{28}\selectfont$A=\pi r^2$}\par\smallskip
{\small base $\pi r$ por altura $r$}
\end{tcolorbox}
\end{minipage}

\vfill
\begin{tcbraster}[raster columns=3,raster equal height=rows,raster column skip=8pt]
\figura{Sector circular}{\begin{tikzpicture}[scale=0.62]\filldraw[fill=fmA!25,draw=fmA,line width=1.3pt] (0,0) -- (0:1.6) arc (0:70:1.6) -- cycle;\draw[fmmuted] (0,0) circle (1.6);\node[font=\scriptsize] at (35:0.55) {$n^\circ$};\end{tikzpicture}}{\large$A=\dfrac{\pi r^2\cdot n}{360}$}
\figura{Corona circular}{\begin{tikzpicture}[scale=0.62]\filldraw[fill=fmB!25,draw=fmB,line width=1.3pt,even odd rule] (0,0) circle (1.6) (0,0) circle (0.9);\draw[fmnavy] (0,0) -- (1.6,0) node[pos=0.75,above,font=\scriptsize]{$R$};\draw[fmnavy] (0,0) -- (0,0.9) node[midway,left,font=\scriptsize]{$r$};\end{tikzpicture}}{\large$A=\pi\,(R^2-r^2)$}
\figura{Trapecio circular}{\begin{tikzpicture}[scale=0.62]\filldraw[fill=fmK!25,draw=fmK,line width=1.3pt] (0:0.9) -- (0:1.6) arc (0:70:1.6) -- (70:0.9) arc (70:0:0.9) -- cycle;\draw[fmmuted] (0,0) circle (1.6) (0,0) circle (0.9);\end{tikzpicture}}{\large$A=\dfrac{\pi\,(R^2-r^2)\cdot n}{360}$}
\end{tcbraster}

\vfill
\begin{exemple}[Ejemplos]
Pizza de 30 cm de diámetro ($r=15$): \ $A=\pi\cdot15^2\approx\boxed{706{,}86\text{ cm}^2}$\\[5pt]
Sector de $60^\circ$ con $r=12$ cm: \ $A=\dfrac{\pi\cdot144\cdot60}{360}=24\pi\approx\boxed{75{,}40\text{ cm}^2}$\\[5pt]
Corona de radios 10 cm y 6 cm: \ $A=\pi\,(100-36)=64\pi\approx\boxed{201{,}06\text{ cm}^2}$\\[5pt]
Círculo de 50,27 cm$^2$: \ $r=\sqrt{50{,}27:\pi}\approx\boxed{4\text{ cm}}$
\end{exemple}

\vfill
\begin{exemple}[Ejemplo: \ ¿qué tiene más pizza, una de 40 cm o dos de 28 cm?]
Una grande: \ $\pi\cdot20^2=400\pi\approx1256{,}6$ cm$^2$ \qquad Dos pequeñas: \ $2\cdot\pi\cdot14^2=392\pi\approx1231{,}5$ cm$^2$\\[3pt]
\textbf{La grande}, aunque $40<56$: el área depende del \textbf{cuadrado} del radio.
\end{exemple}

\vfill
\begin{compte}[Radio, no diámetro]
Las fórmulas usan el \textbf{radio}. Si te dan el diámetro, divídelo entre 2 antes de nada.
\end{compte}

\vfill

\vfill
% ============================================================ 8.9
\clearpage
\seccio{sec9}{Figuras compuestas y problemas}{Sumamos trozos conocidos, o restamos lo que sobra.}

\vfill
\begin{tcbraster}[raster columns=2,raster equal height,raster column skip=10pt]
\begin{tcolorbox}[fmbase,colback=fmK!8,colframe=fmK,coltitle=white,title={\faPlus\ \ Descomponemos y sumamos}]
\begin{minipage}[c]{0.45\linewidth}\centering
\begin{tikzpicture}[scale=0.3]
  \filldraw[fill=fmK!20,draw=fmK!80!black,line width=1.3pt] (0,0) -- (10,0) -- (10,3) -- (4,3) -- (4,8) -- (0,8) -- cycle;
  \draw[fmK!80!black,dashed] (4,0) -- (4,3);
  \node[below,font=\scriptsize] at (5,0) {10}; \node[left,font=\scriptsize] at (0,4) {8}; \node[above,font=\scriptsize] at (2,8) {4}; \node[right,font=\scriptsize] at (10,1.5) {3};
\end{tikzpicture}
\end{minipage}\hfill
\begin{minipage}[c]{0.5\linewidth}
$4\cdot8+6\cdot3$\\ $=32+18=\boxed{50\text{ m}^2}$
\end{minipage}
\end{tcolorbox}
\begin{tcolorbox}[fmbase,colback=fmR!6,colframe=fmR,coltitle=white,title={\faMinus\ \ Figura grande menos trozos}]
\begin{minipage}[c]{0.45\linewidth}\centering
\begin{tikzpicture}[scale=0.3]
  \filldraw[fill=fmR!25,draw=fmR,line width=1.3pt] (0,0) rectangle (8,8);
  \filldraw[fill=white,draw=fmR,line width=1.1pt] (4,4) circle (4);
  \node[below,font=\scriptsize] at (4,0) {10};
\end{tikzpicture}
\end{minipage}\hfill
\begin{minipage}[c]{0.5\linewidth}
$10^2-\pi\cdot5^2=$\\[2pt] $=100-25\pi$\\[2pt] $\approx\boxed{21{,}46\text{ cm}^2}$
\end{minipage}
\end{tcolorbox}
\end{tcbraster}

\vfill
\begin{compte}[El perímetro es solo el contorno exterior]
Las líneas que trazamos por dentro para descomponer la figura \textbf{no cuentan}. \ En la L: \ $P=10+3+6+5+4+8=\boxed{36\text{ m}}$
\end{compte}

\vfill
\begin{exemple}[Ejemplo: \ una ventana rectangular de 1,2 m $\times$ 1,5 m acabada en un semicírculo]
\begin{minipage}[c]{0.22\linewidth}\centering
\begin{tikzpicture}[scale=0.85]
  \filldraw[fill=fmA!15,draw=fmex,line width=1.4pt] (0,0) -- (1.2,0) -- (1.2,1.5) arc (0:180:0.6) -- (0,0);
  \node[below,font=\scriptsize] at (0.6,0) {1,2}; \node[left,font=\scriptsize] at (0,0.75) {1,5};
\end{tikzpicture}
\end{minipage}\hfill
\begin{minipage}[c]{0.75\linewidth}
$A=1{,}2\cdot1{,}5+\dfrac{\pi\cdot0{,}6^2}{2}=1{,}8+0{,}18\pi\approx\boxed{2{,}37\text{ m}^2}$\\[4pt]
Marco: \ $P=1{,}2+2\cdot1{,}5+\dfrac{2\pi\cdot0{,}6}{2}=4{,}2+0{,}6\pi\approx\boxed{6{,}08\text{ m}}$
\end{minipage}
\end{exemple}

\vfill
\begin{exemple}[Ejemplo: \ un camino de 1,5 m de anchura rodea un jardín de 20 m $\times$ 12 m. ¿Qué superficie tiene el camino?]
\begin{minipage}[c]{0.25\linewidth}\centering
\begin{tikzpicture}[scale=0.11]
  \filldraw[fill=fmB!30,draw=fmB,line width=1pt] (0,0) rectangle (23,15);
  \filldraw[fill=fmK!25,draw=fmK,line width=1pt] (1.5,1.5) rectangle (21.5,13.5);
\end{tikzpicture}
\end{minipage}\hfill
\begin{minipage}[c]{0.72\linewidth}
Rectángulo exterior: $(20+3)\times(12+3)=23\times15$\\[3pt]
$A=23\cdot15-20\cdot12=345-240=\boxed{105\text{ m}^2}$
\end{minipage}
\end{exemple}

\vfill
\begin{exemple}[Ejemplo: \ pantalla de 55 pulgadas en formato 16:9 (1 pulgada $=$ 2,54 cm)]
Diagonal: $55\cdot2{,}54=139{,}7$ cm. \ Anchura $16k$, altura $9k$: \ $(16k)^2+(9k)^2=139{,}7^2\ \to\ 337k^2\approx19\,516\ \to\ k\approx7{,}61$\\[3pt]
\boxed{121{,}8\text{ cm}\times68{,}5\text{ cm}}
\end{exemple}

\vfill

\vfill
% ============================================================ 8.10
\clearpage
\seccio{sec10}{El teorema de Pitágoras}{En un triángulo rectángulo, los catetos y la hipotenusa están relacionados.}

\vfill
\begin{minipage}[c]{0.44\linewidth}\centering
\begin{tikzpicture}[scale=0.55]
  \coordinate (C) at (0,0); \coordinate (A) at (4,0); \coordinate (B) at (0,3);
  \filldraw[fill=fmA!25,draw=fmA,line width=1.1pt] (0,0) rectangle (4,-4);
  \filldraw[fill=fmB!25,draw=fmB,line width=1.1pt] (0,0) rectangle (-3,3);
  \filldraw[fill=fmK!22,draw=fmK,line width=1.1pt] (4,0) -- (0,3) -- (3,7) -- (7,4) -- cycle;
  \filldraw[fill=white,draw=fmnavy,line width=1.6pt] (C) -- (A) -- (B) -- cycle;
  \draw[fmnavy,line width=0.8pt] (0,0.4) -- (0.4,0.4) -- (0.4,0);
  \node[font=\large\heavy,text=fmA] at (2,-2) {$a^2=16$};
  \node[font=\large\heavy,text=fmB] at (-1.5,1.5) {$b^2=9$};
  \node[font=\large\heavy,text=fmK!80!black] at (3.5,3.5) {$c^2=25$};
  \node[below,font=\small,text=fmA] at (2,0) {};
  \node[font=\small\heavy,text=fmA,fill=white,inner sep=1pt] at (2,0.35) {$a$};
  \node[font=\small\heavy,text=fmB,fill=white,inner sep=1pt] at (0.35,1.5) {$b$};
  \node[font=\small\heavy,text=fmK!80!black,fill=white,inner sep=1pt] at (2.15,1.75) {$c$};
\end{tikzpicture}
\end{minipage}\hfill
\begin{minipage}[c]{0.53\linewidth}
\begin{idea}[Teorema de Pitágoras]
En un triángulo \textbf{rectángulo}, los lados del ángulo recto son los \cA{\textbf{catetos}} ($a$ y $b$) y el lado opuesto es la \cK{\textbf{hipotenusa}} ($c$), siempre el más largo.\\[4pt]
{\centering\fontsize{22}{26}\selectfont$\cK{c}^2=\cA{a}^2+\cB{b}^2$\par}\smallskip
El área del cuadrado de la hipotenusa es la suma de las áreas de los cuadrados de los catetos: \ $25=16+9$.
\end{idea}
\end{minipage}

\vfill
\begin{tcbraster}[raster columns=2,raster equal height,raster column skip=10pt]
\begin{tcolorbox}[fmbase,colback=fmK!8,colframe=fmK,coltitle=white,title={Queremos la hipotenusa: \textbf{sumamos}}]
{\centering\Large$c=\sqrt{a^2+b^2}$\par}\smallskip
$a=6,\ b=8$: \ $c=\sqrt{36+64}=\sqrt{100}=\boxed{10}$
\end{tcolorbox}
\begin{tcolorbox}[fmbase,colback=fmA!8,colframe=fmA,coltitle=white,title={Queremos un cateto: \textbf{restamos}}]
{\centering\Large$a=\sqrt{c^2-b^2}$\par}\smallskip
$c=13,\ b=5$: \ $a=\sqrt{169-25}=\sqrt{144}=\boxed{12}$
\end{tcolorbox}
\end{tcbraster}

\vfill
\begin{minipage}[t]{0.5\linewidth}
\begin{truc}[Ternas pitagóricas]
Números naturales con $a^2+b^2=c^2$:\\[2pt]
$(3,4,5)\ \ (5,12,13)\ \ (8,15,17)\ \ (7,24,25)$\\[2pt]
Si las multiplicamos, salen más: $(6,8,10)$, $(9,12,15)$\dots
\end{truc}
\end{minipage}\hfill
\begin{minipage}[t]{0.47\linewidth}
\begin{compte}[Antes de restar]
Comprueba que la \textbf{hipotenusa} es el lado más largo. Si la raíz no es exacta, déjala indicada ($\sqrt{74}$) o redondea: $\approx8{,}60$.
\end{compte}
\end{minipage}

\vfill
\begin{exemple}[Ejemplo: \ una escalera de 5 m apoyada en una pared, con el pie a 1,4 m. ¿A qué altura llega?]
\begin{minipage}[c]{0.25\linewidth}\centering
\begin{tikzpicture}[scale=0.42]
  \fill[fmmuted!25] (-0.4,0) rectangle (0,5.2); \draw[fmmuted,line width=1pt] (0,0) -- (0,5.2); \draw[fmmuted,line width=1pt] (0,0) -- (2.6,0);
  \draw[fmB,line width=2pt] (1.4,0) -- (0,4.8);
  \draw[fmnavy,line width=0.7pt] (0,0.35) -- (0.35,0.35) -- (0.35,0);
  \node[below,font=\scriptsize] at (0.7,0) {1,4}; \node[right,font=\scriptsize,text=fmB] at (0.8,2.6) {5}; \node[left,font=\scriptsize\heavy] at (-0.4,2.4) {$h$};
\end{tikzpicture}
\end{minipage}\hfill
\begin{minipage}[c]{0.72\linewidth}
La escalera es la \textbf{hipotenusa} y la altura, un \textbf{cateto}:\\[4pt]
$h=\sqrt{5^2-1{,}4^2}=\sqrt{25-1{,}96}=\sqrt{23{,}04}=\boxed{4{,}8\text{ m}}$
\end{minipage}
\end{exemple}

\vfill
\begin{exemple}[Ejemplo: \ un cateto mide 7 cm más que el otro y la hipotenusa, 13 cm]
$x^2+(x+7)^2=13^2\ \to\ x^2+7x-60=0\ \to\ x=\dfrac{-7\pm17}{2}=5$ \quad \nota{$-12$ no tiene sentido} \quad Catetos: \boxed{5\text{ cm y }12\text{ cm}}
\end{exemple}

\vfill

\vfill
% ============================================================ 8.11
\clearpage
\seccio{sec11}{Pitágoras para calcular áreas y perímetros}{Busca el triángulo rectángulo escondido.}

\vfill
\begin{tcbraster}[raster columns=3,raster equal height=rows,raster column skip=8pt,raster row skip=8pt]
\figura{Diagonal del rectángulo}{\begin{tikzpicture}[scale=0.42]
  \draw[fmnavy,line width=1.3pt] (0,0) rectangle (5,3); \fill[fmK!25] (0,0) -- (5,0) -- (5,3) -- cycle; \draw[fmK,line width=1.6pt] (0,0) -- (5,3);
  \node[below,font=\scriptsize] at (2.5,0) {$b$}; \node[right,font=\scriptsize] at (5,1.5) {$h$}; \node[above left,font=\scriptsize\heavy,text=fmK!80!black] at (2.5,1.5) {$d$};
\end{tikzpicture}}{$d=\sqrt{b^2+h^2}$}
\figura{Diagonal del cuadrado}{\begin{tikzpicture}[scale=0.42]
  \draw[fmnavy,line width=1.3pt] (0,0) rectangle (3.2,3.2); \fill[fmK!25] (0,0) -- (3.2,0) -- (3.2,3.2) -- cycle; \draw[fmK,line width=1.6pt] (0,0) -- (3.2,3.2);
  \node[below,font=\scriptsize] at (1.6,0) {$l$}; \node[right,font=\scriptsize] at (3.2,1.6) {$l$};
\end{tikzpicture}}{$d=\sqrt{l^2+l^2}=l\sqrt2$}
\figura{Altura del isósceles}{\begin{tikzpicture}[scale=0.42]
  \draw[fmnavy,line width=1.3pt] (0,0) -- (4,0) -- (2,3.2) -- cycle; \fill[fmK!25] (2,0) -- (4,0) -- (2,3.2) -- cycle; \draw[fmK,line width=1.6pt] (2,0) -- (2,3.2);
  \node[below,font=\scriptsize] at (3,0) {$\frac b2$}; \node[right,font=\scriptsize] at (3,1.7) {$l$}; \node[left,font=\scriptsize\heavy,text=fmK!80!black] at (2,1.4) {$h$};
\end{tikzpicture}}{$h=\sqrt{l^2-\left(\tfrac b2\right)^2}$}
\figura{Altura del equilátero}{\begin{tikzpicture}[scale=0.42]
  \draw[fmnavy,line width=1.3pt] (0,0) -- (4,0) -- (2,3.46) -- cycle; \fill[fmK!25] (2,0) -- (4,0) -- (2,3.46) -- cycle; \draw[fmK,line width=1.6pt] (2,0) -- (2,3.46);
  \node[below,font=\scriptsize] at (3,0) {$\frac l2$}; \node[right,font=\scriptsize] at (3,1.8) {$l$};
\end{tikzpicture}}{$h=\sqrt{l^2-\left(\tfrac l2\right)^2}\approx0{,}87\,l$}
\figura{Lado del rombo}{\begin{tikzpicture}[scale=0.42]
  \draw[fmnavy,line width=1.3pt] (0,1.6) -- (2.6,0) -- (5.2,1.6) -- (2.6,3.2) -- cycle; \fill[fmK!25] (2.6,1.6) -- (5.2,1.6) -- (2.6,3.2) -- cycle;
  \draw[fmmuted,dashed] (0,1.6) -- (5.2,1.6) (2.6,0) -- (2.6,3.2); \draw[fmK,line width=1.6pt] (5.2,1.6) -- (2.6,3.2);
  \node[below,font=\scriptsize] at (3.9,1.6) {$\frac D2$}; \node[left,font=\scriptsize] at (2.6,2.4) {$\frac d2$};
\end{tikzpicture}}{$l=\sqrt{\left(\tfrac D2\right)^2+\left(\tfrac d2\right)^2}$}
\figura{Altura del trapecio}{\begin{tikzpicture}[scale=0.42]
  \draw[fmnavy,line width=1.3pt] (0,0) -- (6,0) -- (4.6,2.6) -- (1.4,2.6) -- cycle; \fill[fmK!25] (0,0) -- (1.4,0) -- (1.4,2.6) -- cycle; \draw[fmK,line width=1.6pt] (1.4,0) -- (1.4,2.6);
  \node[below,font=\scriptsize] at (0.7,0) {$\frac{B-b}2$}; \node[left,font=\scriptsize] at (0.6,1.5) {$c$};
\end{tikzpicture}}{$h=\sqrt{c^2-\left(\tfrac{B-b}2\right)^2}$}
\end{tcbraster}
{\comic\color{fmmuted}\centering Siempre igual: encuentra el triángulo rectángulo, decide qué lado es la hipotenusa y aplica Pitágoras.\par}

\vfill
\subsection{Distancia entre dos puntos del plano}
\begin{minipage}[c]{0.55\linewidth}
El segmento $AB$ es la \textbf{hipotenusa} de un triángulo rectángulo de catetos $|x_2-x_1|$ y $|y_2-y_1|$:\par\smallskip
{\centering\Large$d(A,B)=\sqrt{(x_2-x_1)^2+(y_2-y_1)^2}$\par}\medskip
$A(1,\,2)$, $B(7,\,10)$: \ $d=\sqrt{6^2+8^2}=\sqrt{100}=\boxed{10}$
\end{minipage}\hfill
\begin{minipage}[c]{0.4\linewidth}\centering
\begin{tikzpicture}[scale=0.34]
  \eixos{-1}{8}{-1}{11}
  \fill[fmK!20] (1,2) -- (7,2) -- (7,10) -- cycle;
  \draw[fmK,line width=1.6pt] (1,2) -- (7,10); \draw[fmA,line width=1.3pt] (1,2) -- (7,2); \draw[fmB,line width=1.3pt] (7,2) -- (7,10);
  \fill[fmnavy] (1,2) circle (5pt) node[left,font=\small] {$A$}; \fill[fmnavy] (7,10) circle (5pt) node[above,font=\small] {$B$};
  \node[below,font=\small,text=fmA] at (4,2) {6}; \node[right,font=\small,text=fmB] at (7,6) {8};
\end{tikzpicture}
\end{minipage}

\vfill
\begin{exemple}[Ejemplos]
\begin{minipage}[t]{0.5\linewidth}
Isósceles de lados 10 y base 12:\\ $h=\sqrt{10^2-6^2}=\sqrt{64}=\boxed{8\text{ cm}}$\\[4pt]
Rombo de diagonales 16 y 12: \ $l=\sqrt{8^2+6^2}=\boxed{10\text{ cm}}$
\end{minipage}\hfill
\begin{minipage}[t]{0.47\linewidth}
Trapecio isósceles de bases 20 y 8, lados 10:\\ $\frac{20-8}{2}=6$, \ $h=\sqrt{10^2-6^2}=\boxed{8\text{ cm}}$\\[4pt]
Equilátero de 6 cm: \ $h=\sqrt{36-9}\approx\boxed{5{,}20\text{ cm}}$
\end{minipage}
\end{exemple}

\vfill

\vfill
% ============================================================ resum
\resumtitol
\begin{tcbraster}[raster columns=2,raster equal height=rows,raster column skip=14pt,raster row skip=20pt,raster before skip=12pt]
\rtile{8.1}{Unidades}{Longitud: de 10 en 10.\\[1ex] Superficie: de 100 en 100.\\[1ex] 1 ha $=10\,000$ m$^2$.}
\rtile{8.2}{Perímetro}{La suma de los lados.\\[1ex] Polígono regular: $n\cdot l$.}
\rtile{8.3}{Rectángulo y cuadrado}{$A=b\cdot h$, \ $A=l^2$.}
\rtile{8.4}{Romboide y triángulo}{$A=b\cdot h$, \ $A=\frac{b\cdot h}{2}$.\\[1ex] La altura, perpendicular.}
\rtile{8.5}{Rombo y trapecio}{$A=\frac{D\cdot d}{2}$, \ $A=\frac{(B+b)\,h}{2}$.}
\rtile{8.6}{Polígonos regulares}{$A=\frac{P\cdot a}{2}$.}
\rtile{8.7}{Circunferencia}{$L=2\pi r$.}
\rtile{8.8}{Círculo}{$A=\pi r^2$.\\[1ex] Sector: $\frac{\pi r^2 n}{360}$.}
\rtile{8.9}{Compuestas}{Sumamos trozos o restamos lo que sobra.}
\rtile{8.10}{Pitágoras}{$c^2=a^2+b^2$.\\[1ex] Hipotenusa: sumamos; cateto: restamos.}
\rtile{8.11}{Aplicaciones}{Busca el triángulo rectángulo escondido.}
\end{tcbraster}
\end{document}
